% ===================================================================
%  CUADRO N.º 11 — La Ciudad y sus monumentos. — El Incendio.
%  Traduction 2026 d'apres texto/io, controlee sur texto/fr et
%  texto/en. Consigne : voir texto/es/10-cuadro-01.tex.
% ===================================================================

%%K t11-apar-1 apar
\VUsaut{2.0mm}
\VUcentre{13.2pt}{90}{TERCERA SERIE}
\VUsaut{3.0mm}
\VUfilet{16mm}
\VUsaut{5.6mm}
\VUtitre{15.0pt}{60}{CUADRO N.º 11}
\VUsaut{1.4mm}
\VUpk{11.4pt}{40}{La Ciudad y sus monumentos. --- El Incendio.}
\VUsaut{2.2mm}

%%K t11-tit-1 sub
\VUpk{10.2pt}{40}{Las afueras.}
\VUsaut{3.0mm}

%%K t11-c1-01-1 p
1. --- Vivo en una gran ciudad situada a cien kilómetros del océano y que,
sin embargo, es un \VUgras{puerto} \textsuperscript{(1)} de primer orden. El río que
la baña tiene cuatrocientos metros de ancho y forma una rada muy hermosa.

%%K t11-c1-02-1 p
2. --- Toda la parte de la ciudad que está junto a los \VUgras{muelles}
\textsuperscript{(2)} tiene un aspecto completamente marítimo e industrial y forma
un \VUgras{arrabal} \textsuperscript{(3)} habitado solo por marinos y obreros. Estos
trabajan en las \VUgras{fábricas} \textsuperscript{(4)} y en la \VUgras{fábrica de gas}
\textsuperscript{(5)}, cuya alta \VUgras{chimenea} \textsuperscript{(6)} y cuyo
\VUgras{gasómetro} se ven desde muy lejos.

%%K t11-c1-03-1 p
3. --- En la Edad Media la ciudad estaba fortificada, y todavía se
encuentran algunos restos de sus murallas, en particular la
\VUgras{puerta} \textsuperscript{(8)} del viejo \VUgras{castillo fuerte}
\textsuperscript{(7)}, flanqueado por sus dos \VUgras{torres} \textsuperscript{(9)}, que hoy
sirven de \VUgras{cárcel}.

%%K t11-c1-04-1 p
4. --- En todo este \VUgras{barrio}, de aspecto muy antiguo, un solo
edificio es relativamente moderno: la \VUgras{aduana} \textsuperscript{(10)}. Está
entre el muelle y la pequeña \VUgras{calle} \textsuperscript{(11)} que lleva al viejo
\VUgras{ayuntamiento} \textsuperscript{(12)}. Este último monumento se reconoce
fácilmente por el gigantesco \VUgras{campanario} \textsuperscript{(13)} que domina
toda la ciudad vieja.

%%K t11-c1-05-1 p
5. --- Al otro lado de la calle \VUgras{se alza} un edificio construido en el
mismo estilo que la aduana: es la \VUgras{bolsa} \textsuperscript{(14)}. Una de sus
fachadas da a una placita en la que está el \VUgras{gobierno civil}
\textsuperscript{(15)}. Nuestra ciudad, en efecto, sin ser una capital, es una
importante cabeza de provincia.

%%K t11-tit-2 sub
\VUsaut{5.0mm}
\VUpk{10.2pt}{40}{La Parte alta de la Ciudad.}
\VUsaut{3.4mm}

%%K t11-c1-06-1 p
6. --- La guarnición, bastante numerosa, se aloja en vastos
\VUgras{cuarteles} \textsuperscript{(16)} muy salubres; se extienden hasta la verja
del \VUgras{parque municipal} \textsuperscript{(17)}. El \VUgras{bulevar} \textsuperscript{(19)}
que lleva a ese parque pasa por detrás de la \VUgras{caja de ahorros}
\textsuperscript{(18)} y desemboca en la plaza de la \VUgras{catedral} \textsuperscript{(20)}.

%%K t11-c1-07-1 p
7. --- Todos estos monumentos se escalonan en anfiteatro sobre una colina
donde está también nuestro amplio \VUgras{instituto} \textsuperscript{(21)}.

%%K t11-c1-07-2 p
Si se baja por una vía algo en cuesta que se junta con la Calle Mayor, se
llega al \VUgras{palacio de justicia} \textsuperscript{(22)}, delante del cual se
extiende un agradable \VUgras{jardín público} \textsuperscript{(26)}. Esta
\VUgras{plaza ajardinada} no es muy grande, pero pone en medio de la ciudad
un oasis de verdor y de frescura. Por eso la frecuentan mucho los
habitantes, a quienes les gusta venir a sentarse bajo el toldo del
\VUgras{café} \textsuperscript{(25)} para escuchar a la banda militar que toca los
jueves y los domingos en el \VUgras{quiosco} \textsuperscript{(27)} puesto entre los
céspedes y los macizos.

%%K t11-c1-08-1 p
8. --- En uno de esos céspedes se alza una \VUgras{estatua} \textsuperscript{(28)} de
bronce, monumento erigido a la memoria de uno de nuestros paisanos que
prestó grandes servicios a su ciudad natal.

%%K t11-tit-3 sub
\VUsaut{4.0mm}
\VUpk{10.2pt}{40}{Los transportes.}
\VUsaut{3.4mm}

%%K t11-c1-09-1 p
9. --- Nuestra ciudad todavía no está alumbrada con luz eléctrica; solo
tiene \VUgras{faroles de gas} \textsuperscript{(29)}. Pero la \VUgras{prensa local}
hace campaña para que ese sistema de alumbrado se mejore, \VUgras{igual que}
ya se ha perfeccionado el \VUgras{sistema de tracción} de los tranvías.
\VUgras{Los} viejos coches \VUgras{tirados} por \VUgras{caballos} han sido
sustituidos por prácticos \VUgras{tranvías eléctricos} \textsuperscript{(31)}, que
llevan rápidamente a los viajeros de un extremo de \VUgras{la ciudad al
otro} por un precio muy \VUgras{barato}.

%%K t11-c1-10-1 p
10. --- Junto al palacio de justicia está el \VUgras{banco} \textsuperscript{(32)},
donde se descuentan los valores comerciales. Los \VUgras{ordenanzas del
banco} \textsuperscript{(33)} van a domicilio a hacer los cobros.

%%K t11-tit-4 sub
\VUsaut{4.0mm}
\VUpk{10.2pt}{40}{Arte e Instrucción.}
\VUsaut{3.4mm}

%%K t11-c1-11-1 p
11. --- El hermoso edificio que hay al lado es el \VUgras{museo}. Contiene a
la vez la \VUgras{biblioteca} \textsuperscript{(34)} de la ciudad, hermosas
\VUgras{colecciones de ciencias naturales} \textsuperscript{(35)} (museo de historia
natural) y una graciosa \VUgras{galería de pintura} \textsuperscript{(36)} que
constituye una espléndida \VUgras{exposición} permanente.

%%K t11-c1-11-2 p
Está abierto al público dos veces por semana, pero los forasteros pueden
visitarlo cualquier día dando una propina al \VUgras{guarda}
\textsuperscript{(37)}. Los \VUgras{pintores} \textsuperscript{(38)} vienen allí a trabajar. Se
los ve \VUgras{delante} de su caballete, con la \VUgras{paleta} en la mano,
copiando los cuadros célebres.

%%K t11-c1-12-1 p
12. --- En la altura, detrás del museo, se distingue la \VUgras{cúpula}
\textsuperscript{(40)} del \VUgras{hospital} \textsuperscript{(39)} y los edificios de la
\VUgras{escuela normal} \textsuperscript{(41)} en la que se formaban los maestros de
nuestras escuelas municipales. En la plaza del Teatro hay una
\VUgras{oficina de correos} \textsuperscript{(43)} instalada en una \VUgras{casa
particular} \textsuperscript{(42)}.

%%K t11-tit-5 sub
\VUsaut{5.0mm}
\VUpk{11.4pt}{40}{El Incendio.}
\VUsaut{3.0mm}

%%K t11-tit-6 sub
\VUpk{10.2pt}{40}{El grito de «¡Fuego!».}
\VUsaut{3.4mm}

%%K t11-c2-01-1 p
1. --- Una noticia alarmante se extiende por la ciudad: ¡acaba de
declararse un incendio en el teatro! Cuando llego a la \VUgras{plaza}
\textsuperscript{(96)}, la multitud ya se apiña en la \VUgras{acera} \textsuperscript{(46)} y
en la \VUgras{calzada} \textsuperscript{(47)}. Los \VUgras{empleados de correos}
\textsuperscript{(44)} y los \VUgras{carteros} \textsuperscript{(45)} salen de la oficina. Un
\VUgras{conductor de tranvía} \textsuperscript{(48)} ha tenido que parar su tranvía
para dejar pasar una \VUgras{ambulancia} \textsuperscript{(50)} municipal que llega
con un \VUgras{cirujano} \textsuperscript{(52)} y un \VUgras{enfermero} \textsuperscript{(51)}
para atender a un \VUgras{herido} \textsuperscript{(53)} que han traído en una
\VUgras{camilla} \textsuperscript{(54)} hasta cerca del \VUgras{quiosco de periódicos}
\textsuperscript{(55)}.

%%K t11-c2-02-1 p
2. --- Una compañía de infantería que volvía del ejercicio se ha parado
para mantener el orden con los \VUgras{guardias municipales} \textsuperscript{(61)} a
caballo. Los \VUgras{jinetes}, cuyos caballos se encabritan, hacen retroceder
a los \VUgras{curiosos} \textsuperscript{(56)} mejor que los \VUgras{soldados}
\textsuperscript{(57)} o los \VUgras{gendarmes} \textsuperscript{(63)}. Por lo demás, estos
últimos están muy ocupados. Uno de ellos indica a los \VUgras{agentes de
policía} \textsuperscript{(64)} adónde hay que llevar a otra \VUgras{víctima}
\textsuperscript{(65)}, un valiente bombero que se ha caído de una escalera.

%%K t11-c2-03-1 p
3. --- El \VUgras{oficial} \textsuperscript{(66)} precisa el sitio donde debe colocarse
la \VUgras{bomba de vapor} \textsuperscript{(67)} que va llegando. Mientras espera esa
potente máquina, ha hecho montar una \VUgras{bomba de mano} \textsuperscript{(68)},
cuya larga \VUgras{manguera} \textsuperscript{(69)} lanza torrentes de agua sobre el
fuego. Seis bomberos la manejan mientras su compañero, que sostiene la
\VUgras{lanza} \textsuperscript{(70)}, dirige el \VUgras{chorro} \textsuperscript{(71)} al centro
del incendio.

%%K t11-c2-04-1 p
4. --- Al otro lado del teatro se ha apoyado la gran \VUgras{escalera}
\textsuperscript{(72)}. Permite a los bomberos acercarse a las llamas que brotan
entre remolinos de \VUgras{humo} \textsuperscript{(75)} negro.

%%K t11-tit-7 sub
\VUsaut{5.0mm}
\VUpk{10.2pt}{40}{Actos de valor.}
\VUsaut{3.4mm}

%%K t11-c2-05-1 p
5. --- El \VUgras{incendio} \textsuperscript{(74)} nació en el vestíbulo del
\VUgras{teatro} \textsuperscript{(73)}, mientras se ensayaba la \VUgras{ópera} cuya
primera \VUgras{representación} debía tener lugar mañana. Por suerte solo
había unos pocos \VUgras{espectadores} \textsuperscript{(76)} en la sala. Los pobres se
han refugiado en el \VUgras{balcón} \textsuperscript{(77)}, adonde valientes
\VUgras{bomberos} \textsuperscript{(86)} acuden a socorrerlos con una escalera y
cuerdas.

%%K t11-c2-06-1 p
6. --- Bajo el \VUgras{pórtico} \textsuperscript{(78)}, donde el \VUgras{cartel}
\textsuperscript{(79)} anuncia la representación, se ve huir a los \VUgras{músicos}
\textsuperscript{(81)} de la \VUgras{orquesta} \textsuperscript{(80)} llevándose sus instrumentos:
\VUgras{violín} \textsuperscript{(81)}, \VUgras{flauta} \textsuperscript{(82)},
\VUgras{violonchelo} \textsuperscript{(83)}, \VUgras{trombón} \textsuperscript{(84)},
\VUgras{tambor} \textsuperscript{(85)}, etc. Se intenta salvar parte del
\VUgras{mobiliario} \textsuperscript{(87)}.

%%K t11-c2-07-1 p
7. --- Los \VUgras{actores} \textsuperscript{(89)} y las \VUgras{actrices}
\textsuperscript{(88)}, \VUgras{cantantes} ellos y ellas, han tenido que huir con su
\VUgras{traje} \textsuperscript{(90)} de teatro. El tenor explica a un
\VUgras{periodista} \textsuperscript{(92)} cómo ha ocurrido. \VUgras{El reportero}
acudió corriendo desde el primer momento con las autoridades: el
\VUgras{gobernador} \textsuperscript{(91)}, el \VUgras{alcalde} \textsuperscript{(93)}, el
\VUgras{comisario de policía} \textsuperscript{(94)} y un \VUgras{juez} \textsuperscript{(95)},
que abrirán una investigación para determinar las responsabilidades.
\VUsaut{6.0mm}
\VUfilet{22mm}
